\section{Conclusions}
\label{sec:conclusion}

In this work, we show that teams of LLM agents can autonomously exploit zero-day
vulnerabilities, resolving an open question posed by prior work
\cite{fang2024llm2}. Our findings suggest that cybersecurity, on both the
offensive and defensive side, will increase in pace. Now, black-hat actors can
use AI agents to hack websites. On the other hand, penetration testers can use
AI agents to aid in more frequent penetration testing. It is unclear whether AI
agents will aid cybersecurity offense or defense more and we hope that future
work addresses this question. Beyond the immediate impact of our work, we hope
that our work inspires frontier LLM providers to think carefully about their
deployments.




\section{Limitations, Ethical Considerations}


Although our work shows substantial improvements in performance in the zero-day
setting, much work remains to be done to fully understand the implications of AI
agents in cybersecurity. For example, we focused on web, open-source
vulnerabilities, which may result in a biased sample of vulnerabilities. We hope
that future work addresses this problem more thoroughly.

A major consideration when conducting research in potentially harmful uses of
LLMs is that malicious actors can use the ideas for nefarious purposes. To help
alleviate such issues, we have elected not to release our code or prompts
publicly as OpenAI has requested that we keep our agents confidential. This is
in line with prior work \cite{fang2024llm, fang2024llm2} and best practice for
cybersecurity \cite{owasp2024vulnerability}. Furthermore, we have disclosed our
findings to OpenAI as part of their responsible disclosure program.
